\section{Sharp planar interactions and a Q-corner Euler criterion}
\label{sec:planar-overlay-refinements}

Arrival \code{3e07f86cf} supplies report 81 while the preceding native
integration is being completed. Its full delivered manuscript, PDF,
alternative implementation and evidence are preserved by placement
\code{6a92806b1}; the committed archive is subsequently retired.
The report strengthens the output bound, extends the mathematical admission
criterion, and isolates a cheaper Euler-only objective. These results are
preserved independently of implementation performance. Its alternative
planar modules use different certificate schemas from report 80; their
patch is retained for comparison rather than overwriting the maintained
producer/checker pair.

\subsection{The interaction constant improves from nine to one}

For the previous section's polygon $P$, set $A_v=m_v-1$ and $A=\sum_vA_v$.
The earlier segment-pair bound is valid but overcounts crossings within
convex cell diagrams. Two diagrams with $f_1,f_2$ full-dimensional cells
have at most $(f_1-1)(f_2-1)$ proper transverse crossings of open edges.
In general position, Euler's formula gives an overlay cell count
$F=f_1+f_2-1+X$. Each overlay cell is one convex intersection of a pair
of cells, so $F\le f_1f_2$. This proves the crossing bound.
For degeneracies, choose a strict interior witness for every retained
minimum function. Sufficiently small generic rational perturbations
preserve these witnesses and each original proper transverse crossing
in a disjoint neighborhood. They do not introduce more retained functions;
applying the general-position bound proves the same inequality. Collinear
overlaps add only existing endpoints, and crossings at existing vertices
are already counted.

\begin{proposition}[Sharp abstract interaction bound]
The complete ray count satisfies
\[
 R\le b+2A+\sum_{u<v}A_uA_v
       \le b+\frac{A(A+3)}2.
\]
For every finite prescribed list of nonnegative $A_v$, the first bound
is attained by a rational affine-potential system. With normal-sector
matching nullity three, the inherited graph count gives $A\le3k+3$,
so $R\le(9k^2+29k+18)/2$. With one varying group, $R\le7k+6$.
\end{proposition}
\begin{proof}
Count original polygon corners, individual diagram vertices and proper
cross-group crossings, using the preceding bounds. Since
$\sum_vA_v^2\ge A$, the second inequality follows.
For sharpness choose $A_v$ parallel lines per group in pairwise nonparallel
rational directions, with small generic offsets so every cross-group
intersection lies inside a polygon containing zero. They give distinct
boundary endpoints and no triple crossings. If $z=n_v\cdot x$, the forms
$h_j=-jz+\sum_{s=1}^j\tau_s$ realize the successive strips with breaks
at $\tau_j$. The vertex count is exactly the first displayed bound.
These systems homogenize to rational potential cones. Realization of the
quadratic lower bound by actual manifold triangulations is not proved.
\end{proof}

\subsection{Certified forced-zero removal admits more feasible cones}

For $C=\{q\ge0:A_\Sigma q=0\}$, let $J$ be the coordinates positive in
some feasible vector. Removing the other allowed types preserves every
feasible standard vector. Summing one positive witness per surviving
coordinate gives a vector strictly positive on $J$. Small perturbations in
any direction of the rebuilt kernel remain nonnegative there, so that
kernel equals the linear span of its feasible cone. Thus feasible-cone
dimension at most three, rather than raw nullity at most three, suffices
after this exact recompilation.

The following single homogeneous LP identifies $J$:
\[
 \max\mathbf1^Tz:\quad A_\Sigma q=0,\ q\ge0,\quad
                     0\le z\le\mathbf1,\quad z\le q.
\]
Scaling and summing the individual feasibility witnesses makes every
$q_j\ge1$ on $J$, giving optimum $|J|$ and unique optimal auxiliary
$z=\mathbf1_J$. The original $q$ need not be unique. Its dual is
\[
 \min\mathbf1^Tu:\quad u,v\ge0,\quad u+v\ge\mathbf1,\quad
                      A_\Sigma^Ty\ge v.
\]
Strong duality supplies independently checkable support witnesses
\[
 A_\Sigma q=0,\quad q\ge0,\quad q_j\ge1\ (j\in J),\qquad
 w=A_\Sigma^Ty,\quad w_j=0\ (j\in J),\quad w_j\ge1\ (j\notin J).
\]
Indeed for any feasible $x$, $w^Tx=y^TA_\Sigma x=0$ is a sum of
nonnegative terms, forcing every off-$J$ coordinate to vanish; $q$ proves
the opposite inclusion. This is a standard maximal-support LP construction
specialized to the sector, without a literature-priority claim.
Adding $\sum q_j=1$ or coordinate caps would invalidate its scaling proof.
Neither delivered implementation installs this preprocessing. At that
delivery all 188 higher-raw-nullity corpus sectors remained outside the
native planar gate. Section~\ref{sec:feasible-span} later implements support
removal from completed Q lists and admits 22 of them to planar enumeration;
the polynomial LP producer itself remains unimplemented.

\subsection{Euler maxima need only the quadrilateral corners}

\begin{proposition}[All-dimensional canonical convexity]
Let $\lambda$ be a linear functional on standard matching vectors with
$\lambda(L_v)\ge0$ for every vertex link. On any nonempty normalized
quadrilateral section $P$,
\[
 \max_{q\in P}\lambda(\operatorname{lift}(q))
 =\max_{a\in\operatorname{Vert}(P)}\lambda(\operatorname{lift}(a)).
\]
Every corner lift spans a non-link standard extreme ray. For compact finite
three-manifold triangulations this applies to Euler characteristic; a
nonpositive maximum excludes every essential normal disc in that sector.
\end{proposition}
\begin{proof}
Let $U(q)$ be the signed linear triangle-potential matching vector with
zero offsets, including singleton vertex groups. Canonical lifting gives
\[
 \lambda(\operatorname{lift}(q))=\lambda(U(q))-
       \sum_v\lambda(L_v)\min_{c\in G_v}\ell_c(q).
\]
Each minimum is concave, so the expression is convex and positively
homogeneous. A convex combination of corners has value at most their
maximum; the reverse inequality is immediate. A Q-corner is extreme in
any minimum cell containing it; the coordinate-face correspondence makes
its lift standard-extreme. Euler values of compact manifold vertex links
are one or two, meeting the sign hypothesis. An essential normal disc has
no vertex-link component and is canonical, by the classical canonical
surface properties reviewed by Burton~\cite{burton-convert}. It has
nonzero quadrilateral vector and positive Euler value, contradicting a
nonpositive maximum. An empty section contains only link surfaces.
\end{proof}

The maximum is of the normalized value $\chi/\sum q_i$, not the raw Euler
characteristic of primitive integer representatives. Positive rescaling
preserves sign, but not numerical maxima. A positive Q-corner can be a
sphere, projective plane or inessential disc, so the essentiality observer
is still required. An essential standard disc need not itself be a
Q-corner. This criterion neither replaces complete standard enumeration
nor supplies a complete sector family.

A primitive integral vector on an admissible standard extreme ray is
connected. Each component is a nonnegative matching vector in the same
compatible sector, so extremality makes it a positive scalar multiple of
the whole vector. Bezout's identity for the primitive coordinates makes
that scalar an integer; it cannot lie strictly between zero and one.
This justifies skipping nonpositive-Euler primitive extreme rays during
disc discovery. Primitivity without extremality does not give this result.

For nullity at most three, a proposed Euler precheck can decide this sign
in $O(1+t+k^2)$ rational operations after kernel construction. Build the
bounded-coefficient linear Euler functional using one representative per
physical face and edge in $O(t)$; project the relevant potentials to the
fixed-dimensional chart and construct $P$ in $O(k^2)$. Aggregate the
linear Euler term and link coefficients in one source pass, then evaluate
the group minima at its at most $k$ corners in $O(k^2)$. A positive witness
lift costs $O(t+k)$. Dense source replay and compressed disc testing are
separate costs. This faster aggregated implementation is not installed;
the maintained complete Q phase already provides the logically valid
nonpositive-Euler exclusion, with its existing proof and costs.
In dimension $d$, active-constraint enumeration tests at most
$\binom{k}{d-1}$ corner systems, including collapsed sections, giving
$k^{O(d)}$ supplied-sector Euler queries. Logarithmic nullity therefore
allows local quasi-polynomial Euler queries, without a global coverage
or essentiality theorem.

\subsection{Arithmetic models, independent evidence and integration limits}

The deterministic sort-and-merge realization has arithmetic bound
$O(1+k^3+(t+k)R)$ and polynomial bit complexity. Both supplied Python
implementations use exact sets/dictionaries; adversarial hash collisions
are not covered by that sharper cubic bound and can add polynomial work.
The geometric bound is not a claim of cubic worst-case Python bit time.
The inherited cofactor estimate bounds primitive triangle coordinates
by $4^s$ and quadrilaterals by $4^{s-1}$ for actual support $s$: the
triangle incidence block is totally unimodular and each quadrilateral
column has absolute column sum at most four. Thus output coordinates
have $O(k)$ bits, even if elementary pieces are exponentially numerous.

Fresh scratch reproduction passes all thirty relevant tests, checks 304
abstract models with 609 compared rays and 1,409 feasible arrangement
candidates, and passes 96 focused certificate/budget/cancellation controls.
Its Euler audit compares 607 declared sectors using exact normalized
maxima and source standard-ray membership. These fresh results are retained
under \code{data/incoming-eaad-81-*}. The delivered all-48-source Regina
audit and 1,344-test historical run are preserved as delivered evidence;
they are not called fresh native runs here.

The conditional recognition statement requires a certified polynomial-size
diagram exterior, a quasi-polynomially constructed complete family of
low-nullity compatible sectors, and exact polynomial-time essential-disc
testing per emitted ray. Local enumeration then preserves quasi-polynomial
time. The complete-family hypothesis remains the substantive missing
theorem. None of the sharpness, forced-zero, convexity or finite audit
results proves it.
